\question[4]
Erläutere, was es bedeutet, dass eine Datenbank in der 3. Normalform vorliegt.
Gib ein Beispiel an für eine Datenbank, die nicht in der 3. Normalform vorliegt,
und erkläre, wie man sie in die 3. Normalform überführen kann.


\begin{solutionbox}{12cm}
Eine Datenbank befindet sich in der 3. Normalform, wenn sie in der 2. Normalform vorliegt
und kein Nicht-Schlüsselattribut von einem anderen Nicht-Schlüsselattribut abhängt.

\vspace{0.2cm}

\textbf{Beispiel für keine 3. Normalform:}

\begin{center}
\begin{tabular}{|l|l|l|l|}
\hline
MitarbeiterID & Name & AbteilungsID & AbteilungsName \\
\hline
1 & Anna Meier & 10 & Buchhaltung \\
2 & Max Schulz & 20 & IT \\
3 & Lea Weber & 10 & Buchhaltung \\
\hline
\end{tabular}
\end{center}

Hier ist \textit{MitarbeiterID} der Schlüssel. Das Attribut \textit{AbteilungsName} hängt aber nicht direkt vom Schlüssel ab,
sondern von \textit{AbteilungsID}. 

\vspace{0.2cm}

\textbf{Überführung in die 3. Normalform:}

\begin{center}
\begin{tabular}{|l|l|l|}
\hline
MitarbeiterID & Name & AbteilungsID \\
\hline
1 & Anna Meier & 10 \\
2 & Max Schulz & 20 \\
3 & Lea Weber & 10 \\
\hline
\end{tabular}
\end{center}

\begin{center}
\begin{tabular}{|l|l|}
\hline
AbteilungsID & AbteilungsName \\
\hline
10 & Buchhaltung \\
20 & IT \\
\hline
\end{tabular}
\end{center}

Nun hängt jedes Nicht-Schlüsselattribut nur noch direkt vom Schlüssel ab.
Die Tabellen liegen damit in der 3. Normalform vor.

\end{solutionbox}
